% TOP_PROBLEM: {"id": 2765, "title": "Kirby Problem 2.17", "category_id": 7, "category": {"id": 7, "name": "topology", "display_name": "Topology", "description": "Properties preserved under continuous deformations.", "slug": "topology", "order_index": 7, "created_at": "2026-07-31T15:26:25.671Z"}, "status": "open", "problem_number": "KP-2.17", "set_id": 11}
% FIRST_POSTED: 2026-10-01
% ENTRY_KIND: statement_only
% UnsolvedMath ID 2765; source catalogue code KP-2.17
% !TeX program = lualatex
% Definition and sources only; this file is not an LLM solution attempt.
\documentclass[11pt,a4paper]{article}
\usepackage[margin=25mm]{geometry}
\usepackage{fontspec}
\setmainfont{lmroman10-regular.otf}[BoldFont=lmroman10-bold.otf,
 ItalicFont=lmroman10-italic.otf,BoldItalicFont=lmroman10-bolditalic.otf]
\usepackage{amsmath,amssymb,mathtools}
\usepackage{unicode-math}
\setmathfont{latinmodern-math.otf}
\AtBeginDocument{\let\setminus\smallsetminus}
\usepackage{xurl,hyperref}
\hypersetup{colorlinks=true,urlcolor=blue}
\setcounter{secnumdepth}{0}
\setlength{\parindent}{0pt}
\setlength{\parskip}{5pt}
\setlength{\emergencystretch}{3em}
\begin{document}
\section*{Kirby Problem 2.17}\label{2765}
{\small UnsolvedMath ID 2765 · KP-2.17 · Topology · Kirby's Problems in Low-Dimensional Topology}
\subsection{Definitions and mathematical statement}
Let $S$ be a closed oriented surface of genus at least two. A geodesic current
on $S$ is a $\pi_1(S)$-invariant Radon measure on the space of unoriented
geodesics in its hyperbolic universal cover, equivalently on unordered pairs
of distinct boundary points. A nontrivial closed curve $\alpha$, taken up to
free homotopy, defines a counting current $[\alpha]$ by counting its lifts;
iterates count with their multiplicities.

For each marked hyperbolic metric $X$ on $S$, a current $c$ has a length
$\ell_X(c)=i(c,L_X)$. Here $L_X$ is the Liouville current and $i$ is the
continuous extension of geometric intersection number, normalized so that
$i([\alpha],L_X)$ is the length of the geodesic representative of $\alpha$.

Suppose there is a closed curve $\gamma$ for which
\[
 \ell_X(c)=\ell_X([\gamma])\quad\hbox{for every marked hyperbolic metric }X.
\]
Must $c$ be a convex combination of closed-curve currents? Explicitly, must
there be finitely many closed curves $\alpha_1,\ldots,\alpha_r$ and real
numbers $t_j\geq0$ with $\sum_jt_j=1$ and
$c=\sum_jt_j[\alpha_j]$? The conclusion concerns equality as measures, not
only equality of lengths for one metric. The curves need not be simple.
\subsection{Short English statement}
If a current has exactly the same length as one closed curve in every hyperbolic metric, must it be a finite convex mixture of closed curves?
\subsection{Sources}
\begin{itemize}
\item[S1] ULAM.AI and the UnsolvedMath Contributors, supplied problems.json, record 2765 (KP-2.17), Kirby's Problems in Low-Dimensional Topology. Catalogue identity and source transcription; CC BY 4.0.
\url{https://huggingface.co/datasets/ulamai/UnsolvedMath}.
\item[S2] K3: A New Problem List in Low-Dimensional Topology, Problem 2.17. Expanded formulation of the supplied catalogue statement.
\url{https://math.berkeley.edu/sites/default/files/surv-295-ruberman-watermarked-author-pdf.pdf}.
\end{itemize}
\end{document}
